\section{Families with one critical point}\label{sec:morse}

We now establish the obstruction used in the proof. A family is \emph{homogeneous at infinity} if it agrees, outside one fixed ball in every fiber, with a family of degree-two homogeneous functions. Smoothness of a homogeneous function below means smoothness away from the origin, together with its $C^1$ extension at the origin.

\begin{proposition}\label{prop:one-critical}
Let $B=S^9$ and let $p\geq20$. Suppose that $f_b:\R^{2p}\to\R$, $b\in B$, is a smooth family satisfying the following conditions:
\begin{enumerate}
\item it is homogeneous at infinity, with smooth homogeneous data in $T_{p,p}$;
\item each $f_b$ has exactly one critical point, which is nondegenerate and varies smoothly with $b$;
\item the stable spherical class of the family at infinity is zero in $[B,BG]$.
\end{enumerate}
Then its stable tube class is zero in $[B,\T]$.
\end{proposition}

The proof occupies this section and Section~\ref{sec:obstruction}. We first fix the local quadratic model and explain the necessary Morse theory on the noncompact fibers. Comparing a one-critical function with a smooth $h$-cobordism by removing a normal neighborhood of its unstable disk is also central to Kragh's account of generating functions and the Hatcher--Waldhausen map \cite{Kragh2026}. Here the data at infinity vary with the parameter, and we keep their smooth tube class throughout the comparison.

\subsection{Normalization and the negative sphere bundle}

Translate the critical section to the origin by compactly supported fiber diffeomorphisms. Such a family is obtained by flowing a cutoff of the constant translation vector field: choose its plateau to contain all line segments from zero to the critical section. Subtract the critical value using a cutoff equal to one on the ball of radius $A$ and zero outside the ball of radius $2A$, with $A$ sufficiently large. On its transition annulus the derivative of this correction is $O(A^{-1})$, whereas the original gradient has norm at least $cA$. Thus no new critical points appear, and the data at infinity are unchanged.

Let $Q_b$ be the quadratic term at the resulting critical point. Taylor expansion, uniformly over $B$, gives
\[
 f_b(w)-Q_b(w)=O(\norm{w}^3),\qquad
 d_w(f_b-Q_b)=O(\norm{w}^2).
\]
A radial cutoff on a sufficiently small ball replaces $f_b$ by $Q_b$ near zero. The extra derivative is $O(\norm{w}^2)$ on the transition annulus, while $\norm{\nabla Q_b(w)}\geq c_0\norm{w}$ for a uniform $c_0>0$. We have therefore arranged
\begin{equation}\label{eq:local-quadratic}
 f_b(0)=0,\qquad f_b=Q_b\text{ near }0,
\end{equation}
without changing the unique critical point or the infinity data.

\begin{lemma}\label{lem:morse-sublevels}
Suppose a family satisfying~\eqref{eq:local-quadratic} has one nondegenerate critical point and smooth homogeneous data of type $T_{k,l}$ at infinity, where $k,l\geq3$. Its Morse index is $k$. Moreover, the unit sphere bundle of its negative eigenspaces is fiber homotopy equivalent to the family of negative domains at infinity.
\end{lemma}

\begin{proof}
Choose a radius $a>0$ inside the common quadratic neighborhood and a radius $R_0$ with an inner collar on which the family is already homogeneous. Take a small uniform $\delta>0$ and write
\[
 M_-^b=\{f_b\leq-\delta\},\qquad M_+^b=\{f_b\leq\delta\}.
\]
We first show that $M_-^b$ is homotopy equivalent to the negative domain at infinity. On its compact part $\norm{w}\leq R_0$, choose a descending field, radial outward near $\norm{w}=R_0$. This boundary prescription is possible because homogeneity gives
\[
 df_b(\partial_r)=2f_b/r<0
\]
there. Interpolate with negative gradient away from that collar. The conditions are strict linear inequalities, so these choices can be patched smoothly, also over $b$.

The field cannot leave through $f_b=-\delta$. It must reach the outer sphere in finite time: on the compact region its rate of descent has a positive lower bound, while $f_b$ has a finite lower bound. The exit is transverse. Stopping the flow at its first exit and fixing all points outside the ball gives a deformation retraction onto
\[
 M_-^b\cap\{r\geq R_0\}.
\]
The globally stopped homotopy is continuous, which is sufficient here. If $h_b$ denotes the angular function at infinity, angular projection of this outer portion has, over $u$ with $h_b(u)<0$, the radial fiber
\[
 \left[\max\left\{R_0,\sqrt{\frac{\delta}{-h_b(u)}}\right\},\infty\right).
\]
Contracting each such interval to its continuous lower endpoint proves the homotopy equivalence. The open and closed negative domains have the same homotopy type, since zero is a regular value. In particular,
\begin{equation}\label{eq:negative-sublevel-type}
 M_-^b\simeq S^{k-1}.
\end{equation}

On $f_b\geq\delta$, normalized negative gradient has $df_b$-derivative $-1$ and bounded speed. Indeed the gradient is bounded away from zero on the compact part and grows at least linearly at infinity. Extend this to a bounded smooth field by cutting off below $\delta$, away from zero. Each point with $f_b>\delta$ reaches the level $\delta$ after time $f_b-\delta$. Stopped flow gives a deformation retraction of $\R^{k+l}$ onto $M_+^b$, so
\begin{equation}\label{eq:positive-sublevel-type}
 M_+^b\simeq *.
\end{equation}

We give the local handle argument to justify its use for these noncompact sublevels. If the Morse index is $j$, choose local linear coordinates with \mbox{$Q_b=-\norm{x}^2+\norm{y}^2$}, where $x\in\R^j$. The coordinate changes and their inverses have uniformly bounded norms; a global frame is not needed. Choose a smooth function $\mu$ on $[0,\infty)$ with
\[
 -1<\mu'\leq0,\qquad \mu(0)>\delta,
 \qquad \mu(s)=0\text{ for }s\geq2\delta.
\]
Modify the function in this neighborhood to
\[
 \widetilde f_b=f_b-\mu(\norm{x}^2+2\norm{y}^2).
\]
The modification is supported where $f_b<\delta$. Its derivatives in the chart are
\[
 -2(1+\mu')x,\qquad 2(1-2\mu')y.
\]
Thus zero is its only critical point and has value below $-\delta$. We have $\{\widetilde f_b\leq\delta\}=M_+^b$. Bounded normalized descent on the slab between $-\delta$ and $\delta$ retracts this set, relative to $M_-^b$, onto $\{\widetilde f_b\leq-\delta\}$.

The latter set retracts, relative to $M_-^b$, onto
\[
 M_-^b\cup D,\qquad D=\{y=0,\ \norm{x}^2\leq\delta\}.
\]
For membership of $D$, note that $s+\mu(s)\geq\mu(0)>\delta$ for $0\leq s\leq\delta$. At an added point outside $M_-^b$, shrink $y$ radially to zero if $\norm{x}^2\leq\delta$, and otherwise to length $\sqrt{\norm{x}^2-\delta}$. The modified function decreases during this motion, and it agrees continuously with fixing $M_-^b$ at the boundary. This proves the claimed relative deformation.

The relative homology of $(M_+^b,M_-^b)$ is therefore that of one $j$-disk modulo its boundary. Equations~\eqref{eq:negative-sublevel-type}--\eqref{eq:positive-sublevel-type} force $j=k$, and the boundary map makes $\partial D$ a generator of $H_{k-1}(M_-^b;\Z)$. Within the local negative cone, this sphere expands to $aS(E_b^-)$, where $E_b^-$ is the negative eigenspace in the original Euclidean metric.

Flow this sphere to the outer sphere using the descending field above and project angularly. The construction varies smoothly with $b$. On each fiber it takes a homology generator to a generator in the negative domain. Both source and target have the homotopy type of the simply connected sphere $S^{k-1}$. They have CW type, so the integral homology isomorphism is a homotopy equivalence by the homological Whitehead theorem \cite[Corollary~4.33]{Hatcher2002}. The fiberwise criterion for maps of bundles over a CW base makes it a fiber homotopy equivalence. This identifies the claimed spherical fibration.
\end{proof}

Return to $k=l=p$ in Proposition~\ref{prop:one-critical}. By Lemma~\ref{lem:morse-sublevels}, its spherical hypothesis says that the stable class of $E^-$ has zero image under $BJ$. The injectivity of $\pi_9BO\to\pi_9BG$ established in Section~\ref{sec:tubes} implies that this bundle class is zero. The map $\operatorname{Gr}_p(\R^{2p})\to BO(p)$ has homotopy fiber the $(p-1)$-connected Stiefel manifold $V_p(\R^{2p})$, so this map is $p$-connected. The stabilization $BO(p)\to BO$ is also $p$-connected. For $p\geq20$ both induce isomorphisms on $\pi_9$. Hence the negative-eigenspace map is nullhomotopic in this finite Grassmannian.

The space of nondegenerate forms of signature $(p,p)$ retracts onto the corresponding sign forms. We can consequently choose a smooth homotopy $Q_{b,s}$ of nondegenerate forms from $Q_b$ to the fixed $q_{p,p}$, with uniform lower gradient bound $c_1\norm{w}$. Insert this homotopy into a much smaller punctured quadratic ball, using a radial parameter $s(r)$ which is one near zero, zero near the outer boundary, and satisfies
\[
 r\abs{s'(r)}\leq\varepsilon.
\]
Such a cutoff exists on an arbitrarily long logarithmic radius interval. Differentiating $Q_{b,s(r)}(w)$ adds an error $O(\varepsilon r)$ to its quadratic gradient. A sufficiently small $\varepsilon$ excludes new critical points. The parameter is constant near zero, so the resulting function is smooth and equals the fixed $q_{p,p}$ there.

\subsection{A smoothing lemma for stabilization}

Direct sums of homogeneous functions are naturally $C^1$. The following explicit adjustment restores smoothness away from the total origin and will also be used in Hamiltonian transport.

\begin{lemma}\label{lem:stabilization-smoothing}
Let $g_b(w)$ be a smooth homogeneous family with $\norm{\nabla g_b(w)}\geq c\norm{w}$, and let $q'(z)$ be a fixed nondegenerate quadratic form. Put $W=(w,z)$ and $G_0(b,W)=g_b(w)+q'(z)$. Suppose that $F_b(W)$ is smooth and that
\[
 \abs{F_b-G_0}\leq M_0,\qquad
 \norm{d_WF_b-d_WG_0}\leq M_1
\]
uniformly. Then, outside an arbitrarily large fixed ball, $F$ can be modified to agree exactly with a smooth homogeneous family $G_\eta$. The modification creates no new critical points, preserves a neighborhood of every original critical point, and the homogeneous family $G_\eta$ is homotopic through regular $C^1$ homogeneous functions to $G_0$.

The construction is uniform over compact parameter spaces, or over any parameter domain on which the bounds used here are uniform.
\end{lemma}

\begin{proof}
There is $c_0>0$ such that $\norm{\nabla_WG_0}\geq c_0\norm{W}$. For $W\ne0$ choose a smooth cutoff $\rho_\eta$ which vanishes on $[0,\eta]$, equals one on $[2\eta,\infty)$, and has derivative bounded by $C/\eta$, and set
\[
 G_\eta(b,W)=g_b(w)\rho_\eta\left(\frac{\norm{w}}{\norm{W}}\right)+q'(z).
\]
The potentially nonsmooth term vanishes in an open cone around $w=0$. Elsewhere it is smooth, so $G_\eta$ is smooth away from $W=0$. On the support of $1-\rho_\eta$, we have $\norm{w}\leq2\eta\norm{W}$. Homogeneity and the angular derivative bound give
\begin{equation}\label{eq:angular-smoothing-estimates}
 \abs{G_\eta-G_0}\leq C_2\eta^2\norm{W}^2,
 \qquad
 \norm{d_WG_\eta-d_WG_0}\leq C_2\eta\norm{W}.
\end{equation}
In the second estimate the derivative of $\norm{w}/\norm{W}$ costs at most $\norm{W}^{-1}$ where the cutoff changes. For small $\eta$, the straight homotopy from $G_0$ to $G_\eta$ retains a linear nonzero-gradient bound.

Let $\alpha_A$ vanish for $\norm{W}\leq A$, equal one for $\norm{W}\geq2A$, and have derivative bounded by $C_3/A$. Replace $F$ by
\[
 F_A=(1-\alpha_A)F+\alpha_AG_\eta.
\]
On the transition annulus its derivative differs from $d_WG_0$ by at most
\[
 M_1+C_2\eta\norm{W}
 +\frac{C_3}{A}\bigl(M_0+C_2\eta^2\norm{W}^2\bigr).
\]
Relative to $\norm{W}$ this is $O(\eta+\eta^2+A^{-1})$. Choose $\eta$ small and then $A$ large so that the bound is less than $c_0\norm{W}/2$. No critical points occur in the annulus or its exterior. The original critical points already lie in a uniform ball, since $d_WF$ differs from a linearly bounded-below gradient by at most $M_1$. Taking $A$ beyond that ball preserves their neighborhoods. The function is unchanged near zero and thus smooth everywhere.
\end{proof}

Adding positive squares increases the normal disk dimension of the negative
tube $S^{k-1}\times D^l$ while leaving the dimension of its sphere core
$S^{k-1}\times\{0\}$ fixed. We use this freedom to allow the final
deformation of these cores through embeddings in Section~\ref{sec:obstruction}.

Apply the lemma to the normalized $f_b(w)+\norm{z}^2$. Its difference from the homogeneous sum, and that of its vertical derivative, are bounded because $f_b-g_b$ is supported in a fixed $w$-ball. We obtain a family, still with one critical point and with a fixed local model, whose infinity data have type $T_{k,l}$ with
\begin{equation}\label{eq:unbalanced-indices}
 k=p,\qquad l>k+9.
\end{equation}
Both indices remain at least twenty. Its local model is $q_{k,l}$. If its tube family is null at this stage, adding the same number of negative squares as the positive squares just added gives nullity in the diagonal stable system. Fixed reordering of coordinate blocks does not affect nullity. It remains to prove that this enlarged tube family is null, which is the purpose of the next section.
